\hypertarget{examples}{%
\section{Examples}\label{sec:examples}}

The proposed formulations are evaluated with five well-known  examples from the literature.
% The first set originates from the study by \cite{jelenić1999geometrically}, and is used to assess conformance with the requirements %principles
% of strain objectivity and path independence.
% The second set assesses the 
% accuracy against available 
% exact solutions of the rod differential equations.
% The following two examples
% %Section~\ref{sec:helix} and Section~\ref{sec:hockling}
% %investigate
% %conclude with 
% %two demanding examples
% \FCF{are selected}
% %in order
% to demonstrate %which illustrate
% the performance of the proposed formulations under very %extremely
% large rotations.
The simulations
are performed in the Matlab based framework FEDEASLab \citep{filippou2004fedeaslabb}.
The deformed shape renderings
were done in collaboration with \cite{vouvakismanousakis2023opensees_model_generator,perez2024openseesrt}.
All examples use %In all examples 
a simple %Saint Venant-Kirchhoff
hyperelastic material
with %is employed where 
the generalized stress \(\boldsymbol{s}\) and  strain \(\boldsymbol{e}\)
%are related 
related in the reference basis $\{\mathbf{D}, \mathbf{D}_\alpha\}$
by:
\[
\left(\begin{array}{c}
N   \\
V_1 \\
V_2 \\
T   \\
M_1 \\
M_2 \\
\end{array}\right)=\left[\begin{array}{cccccc}
E A & 0 & 0 & 0 & 0 & 0 \\
0 & G {A} & 0 & 0 & 0 & 0 \\
0 & 0 & G {A} & 0 & 0 & 0 \\
0 & 0 & 0 & G J & 0 & 0 \\
0 & 0 & 0 & 0 & E I_{1} & E I_{12} \\
0 & 0 & 0 & 0 & E I_{12} & E I_{2} \\
\end{array}\right]\left(\begin{array}{c}
\epsilon \\
\gamma_1 \\
\gamma_2 \\
\tau \\
\kappa_1 \\
\kappa_2 \\
\end{array}\right)
\]
%and
with 
the constitutive %FCF
parameters \(G\), \(E\), \(I_\alpha\), \(I_{12}\) and \(A\)
representing
standard cross-sectional properties.
A square cross-section is used in the renderings for depicting
%in order to depict
the rotations clearly,
and should not be interpreted as a representation of the cross-section geometry.
The reference configuration is depicted as a wireframe in the renderings, and the deformed configuration is shaded solid.
%The standard Euclidean basis $\{\mathbf{E}_1, \mathbf{E}_2, \mathbf{E}_3\}$
%\FCF{underlies the geometry definition for all problems
% The \FCF{geometry} description 
% of the problems refer to \FCF{uses}
The problems are described in reference to the standard Euclidean basis $\{\mathbf{E}_1, \mathbf{E}_2, \mathbf{E}_3\}$,
which is also depicted in the renderings. 
All simulations use the relative work increment as the convergence criterion. 
%The work increment is used %in all problems to indicate convergence
%as convergence criterion throughout the analyses.
% The wrapped elements under consideration are suceptible to shear locking,
% which is mitigated by the use of reduced Gauss-Legendre quadrature of order \(n-1\), for $n$-noded elements.
The state determination of all wrapped elements uses   
% a uniformly reduced 
Gauss-Legendre quadrature of order \(n-1\), where $n$ is the number of element nodes. %, as is typical for displacement-interpolated beam elements.
The implementation of these simulations is available at \url{https://github.com/claudioperez/GeometricTransformations}.
% Further details of the implementation in FEDEASLab are available in.

